\section{Innledning}
Vibrasjonsmåling er en vanlig metode for å følge med på tilstanden til roterende maskiner. Poenget er å oppdage endringer i vibrasjonsbildet før feilen utvikler seg til havari \parencite{vibrasjon1, tbv1}. Frekvensinnholdet i signalet kan gi informasjon om hvor vibrasjonen kommer fra, for eksempel om den skyldes ubalanse i akselen eller en lokal skade i lageret \parencite{nsen13306}. Vibrasjonsovervåking utnytter P-F-intervallet, altså tidsrommet mellom første detekterbare feil og funksjonssvikt, noe som gir mulighet til å planlegge vedlikehold i tide \parencite{tbv2}.

Formålet med denne laboratorieøvelsen er å måle vibrasjoner på et R10-kulelager montert i en RK4-vibrasjonsrigg, identifisere grunnfrekvensen og de karakteristiske lagerfrekvensene, og vurdere lagerets tilstand ut fra amplitudene i frekvensspekteret.
